\documentclass[10pt,a4paper]{amsart}
\usepackage[margin=19mm]{geometry}
\usepackage{amsmath,amssymb,mathtools}
\usepackage[scr=rsfs,cal=euler]{mathalfa}
\usepackage{xfrac}
\usepackage[unicode,hidelinks,bookmarks=false]{hyperref}
\newcommand{\ie}{i.e.\ }
\newcommand{\cf}{cf.\ }
\newcommand{\resp}{respectively\ }
\newcommand{\Subsection}{\subsection}
\providecommand{\nobreakdash}{\nobreak\hbox{-}}
\setlength{\emergencystretch}{2em}
\multlinegap0pt
\usepackage{bm}
\usepackage{bbm}
\usepackage{dsfont}
\usepackage{xspace}
\usepackage{cancel}
\usepackage{mathtools}
\usepackage{amsthm}
\makeatletter
\@ifundefined{theorem}{\newtheorem{theorem}{Theorem}}{}
\@ifundefined{lemma}{\newtheorem{lemma}[theorem]{Lemma}}{}
\@ifundefined{proposition}{\newtheorem{proposition}[theorem]{Proposition}}{}
\makeatother
\makeatletter
\newcommand{\R}{\mathbb{R}}
\newcommand{\Z}{\mathbb{Z}}
\newcommand{\bigO}{\mathcal{O}}
\newcommand{\dd}{\mathrm{d}}
\expandafter\ifx\csname example\endcsname\relax
\newenvironment{example}
{\pushQED{\qed}\renewcommand{\qedsymbol}{$\diamond$}\examplex}
{\popQED\endexamplex}
\fi
\newcommand{\ol}{\Lambda}
\makeatother
\title[Bivariate exponential integrals and edge-bicolored graphs]{Bivariate exponential integrals and edge-bicolored graphs}
\author{Michael Borinsky, Chiara Meroni, Maximilian Wiesmann}
\date{Source: arXiv:2409.18607v1 (CC BY 4.0)}
\begin{document}
\maketitle

\noindent\emph{Source and licence.} Bivariate exponential integrals and edge-bicolored graphs. Version arXiv:2409.18607v1 (CC BY 4.0), distributed under the Creative Commons Attribution 4.0 International licence, as recorded in the arXiv licence metadata for this exact version and shown on the abstract page https://arxiv.org/abs/2409.18607v1. Full rights notice: licenses/arxiv-2409.18607-CC-BY-4.0.txt.\par\smallskip

\noindent\textit{Editorial scope.} Counted unit: the Proposition whose source label is \texttt{prop:AnAsy} in the article (source0.tex lines 755--775, source label \texttt{prop:AnAsy}), delivered together with the article's own complete original proof of it (source0.tex lines 817--861, one \texttt{proof} environment of 45 lines). No mathematics is written, repaired or re-derived: statement and proof are the article's own text line for line. The proof is detached from the statement in the source: the article states the result at lines 755--775 and prints its proof at lines 817--861, and the proof environment's own header names the statement, so the association is the article's own. Required context retained uncounted: eq:coeffs\_extraction\_gen\_funct (lines 192--196); prop:laplace (lines 211--230); lmm:AnHomo (lines 783--792). Every definition, display and label of those lines is retained unchanged. Omitted from this package and not redistributed: 1--191, 197--210, 231--754, 776--782, 793--816, 862--1256. Nothing inside a retained excerpt is omitted or altered: every retained body is delivered exactly as the article's lines are written, so no omission receipt of any kind accompanies this package. Citations: the retained excerpts carry no citation command at all, so no bibliography entry of the article is transcribed and no reference is redistributed. Reference adaptation: none was needed and none was made; every reference inside the delivered unit resolves inside it, so no unresolved reference or citation is produced. Printed slips kept literal: no printed word, symbol, hypothesis or number of the counted statement or of its proof was altered. Layout and class: the article's own class is \texttt{ourlematemasingle} with packages including hyperref, fontenc, inputenc, anyfontsize, amsmath, amssymb, amsthm, xcolor, paralist, tikz, pgfplots; the wrapper is the lane's pinned \texttt{amsart} editorial wrapper at 10pt on A4, which restates the theorem-like environments the retained text needs and adds the article's own macro definitions that the retained text needs (\texttt{\textbackslash R}, \texttt{\textbackslash Z}, \texttt{\textbackslash bigO}, \texttt{\textbackslash dd}, \texttt{\textbackslash ol}), with extra wrapper packages (bm, bbm, dsfont, xspace, cancel, mathtools). The delivered numbering is the unit's own. Assets: no retained span contains a figure environment, a graphics-inclusion command, a PDF-page inclusion, a table or a TikZ picture; no image file is copied into this package, the package declares no asset dependency and no image file is redistributed with it. Licence: version arXiv:2409.18607v1 of this article is distributed under the Creative Commons Attribution 4.0 International licence, as recorded in the arXiv licence metadata for this exact version and shown on the abstract page https://arxiv.org/abs/2409.18607v1. A full rights notice is delivered at licenses/arxiv-2409.18607-CC-BY-4.0.txt. Characters: every retained excerpt is pure ASCII, so the delivered engine, XeLaTeX with its default Latin Modern text font, reports no missing character.
\begin{equation}
\label{eq:coeffs_extraction_gen_funct}
    \sum_{s,t \geq 0} a_{s,t}(\lambda)\, x^s y^t =
\exp\left(\sum_{\substack{u,w \geq 0\\u+w \geq 1}} \lambda_{u,w} \frac{x^u y^w}{u!w!}\right) \in \mathcal R[[x,y]].
\end{equation}

\begin{proposition}
\label{prop:laplace}
If $I(z)$, $g$, $D$ and the coefficients $\ol_{u,w}$ are related as above, then
\[
I(z) \sim \sum_{n\geq 0} A_n z^{-n},
\]
for large $z$, where $A_n$ is the coefficient of $z^{-n}$ in the formal power series
\[
    \sum_{s,t\geq 0} z^{-(s+t)} (2s-1)!! \cdot (2t-1)!! \cdot a_{2s,2t} (z \cdot \ol) \in \R[[z^{-1}]]\,,
\]
where $(2s-1)!! = (2s-1)(2s-3)\cdots 3\cdot 1$ and $a_{2s,2t} (z \cdot \ol)$ is the polynomial $a_{2s,2t}(\lambda)$ defined in \eqref{eq:coeffs_extraction_gen_funct},
with 
\begin{equation}
    \label{eq:substitution_lambda_Lambda}
    \lambda_{u,w} = \left\{ \begin{array}{ll}
        0 & \text{for}~ u,w\geq 0 ~\text{and}~ 1 \leq u+w < 3, \\
        z\,\Lambda_{u,w} & \text{for}~ u,w\geq 0 ~\text{and}~ u+w \geq 3 .\\
    \end{array} \right.
\end{equation}
\end{proposition}

\begin{lemma}
\label{lmm:AnHomo}
Let $M = \frac{k}{k-2}$ and $K=\frac{2}{k-2}$.
For a given integer $n \geq 0$ such that $nK$ and $nM$ are integers, we have
\begin{gather*}
A_n = \frac{2^{nM} (nM)! }{ 2\pi \cdot (nK)!} 
\int_{-\pi}^\pi  V(\cos\varphi ,\sin \varphi)^{nK} \dd \varphi.
\end{gather*}
If $nK$ or $nM$ is not an integer then $A_n =0$.
\end{lemma}

\begin{proposition}
\label{prop:AnAsy}
Let $M = \frac{k}{k-2}$ and $K=\frac{2}{k-2}$.
If $A_n$, $V$, $\ol_{u,w}$ and $\Phi$ are related as described above and all extrema in $\Phi$ are non-degenerate, then
$$
A_n \sim 
\begin{cases}
\frac{1}{2\sqrt{2}\pi}k^{nM+\frac{1}{2}}K^{n-\frac{1}{2}}\Gamma(n) 
\sum\limits_{(x,y) \in \Phi} \frac{V(x,y)^{nK}}{\sqrt{B(x,y)}}
&\text{ if } nK, nM \in \Z,
\\
0 &\text{ else,}
\end{cases}
$$
where
$$
B(x,y) = 
k^2 - 
\frac{ \frac{\partial^2 V}{\partial x^2}(x,y) + \frac{\partial^2 V}{\partial y^2}(x,y)}{ V(x,y) } \quad \text{ for } (x,y)\in S^1.
$$
\end{proposition}

\begin{proof}[Proof of Proposition~\ref{prop:AnAsy}]
We are interested in the cases in which $A_n\neq 0$.
When $n$ is large, the main contribution to the integral 
in the statement of Lemma~\ref{lmm:AnHomo} comes 
from angles $\varphi$ where $|V(\cos \varphi,\sin \varphi)|$ is 
maximal. Let $\varphi_c$ be the location of such a maximum.
By definition, we have $(\cos \varphi_c,\sin \varphi_c) \in \Phi$.
Near this maximum, we get the Taylor expansion
$$
f_{\varphi_c}(\varphi):=
\log \frac{V(\cos \varphi, \sin \varphi)}{V(\cos \varphi_c, \sin \varphi_c)}
=
-
B(\cos \varphi_c, \sin \varphi_c)
\frac{(\varphi -\varphi_c)^2}{2}
+\bigO((\varphi -\varphi_c)^3),
$$
where $B(\cos \varphi_c, \sin \varphi_c)$ is defined as in the statement.
Because $\varphi_c$ is a maximum of $|V(\cos \varphi_c, \sin \varphi_c)|$, 
we have $B(\cos \varphi_c, \sin \varphi_c) \geq 0$. Our assumption that all the maxima are non-degenerate hence implies that $B(\cos \varphi_c, \sin \varphi_c) > 0$.\par 

We may therefore write, for some sufficiently small $\varepsilon > 0$,
\begin{gather*}
A_n = \frac{2^{nM} (nM)! }{ 2\pi \cdot (nK)!} 
\sum_{(\cos \varphi_c,\sin \varphi_c) \in \Phi}
\!\!\!V(\cos \varphi_c, \sin \varphi_c)^{nK}
\int_{\varphi_c-\varepsilon}^{\varphi_c+\varepsilon} e^{nK f_{\varphi_c}(\varphi)} \dd \varphi + R_1(n,\varepsilon).
\end{gather*}
From the Taylor expansion of the function $f_{\varphi_c}(\varphi)$ and
Lemma~\ref{lmm:AnHomo},
it follows by the same reasoning as in the proof of Proposition~\ref{prop:laplace} that the remainder term respects the bound
$|R_1(n,\varepsilon)| \leq C_1 \exp({-C_2 n \varepsilon^2})$
with some constants $C_1,C_2>0$.
Specifying $\varepsilon = n^{-\gamma}$ with $\gamma \in (\frac13,\frac12)$ allows to truncate the Taylor expansion of $f_{\varphi_c}$ after the second term without changing asymptotic behavior in the $n\rightarrow \infty$ limit. Hence,
$$
\int_{\varphi_c-\varepsilon}^{\varphi_c+\varepsilon} e^{nK f_{\varphi_c}(\varphi)} \dd \varphi =
\int_{-\varepsilon}^{\varepsilon} \exp\left(-nK \frac{B(\cos\varphi_c,\sin \varphi_c)}{2} \varphi^2\right) \dd \varphi + R_2(n, \varepsilon).
$$
The remainder term fulfills $|R_2(n, \varepsilon)| < C_3 n^{\frac12} \varepsilon^3$ for some $C_3>0$.
Again, as in the proof of Proposition~\ref{prop:laplace}, we may complete the Gaussian integral to find that 
$$
\int_{\varphi_c-\varepsilon}^{\varphi_c+\varepsilon} e^{nK f_{\varphi_c}(\varphi)} \dd \varphi = \sqrt{ \frac{2\pi}{nK \cdot B(\cos\varphi_c,\sin \varphi_c)}} + \bigO (n^{-1}).
$$
The result follows from Stirling's formula $\Gamma(n) \sim \sqrt{2\pi n^{-1}} n^n e^{-n}$ as $n \rightarrow \infty$.
\end{proof}

\end{document}
